\documentclass{article}

% if you need to pass options to natbib, use, e.g.:
%     \PassOptionsToPackage{numbers, compress}{natbib}
% before loading neurips_2026

% The authors should use one of these tracks.
% Before accepting by the NeurIPS conference, select one of the options below.
% 0. "default" for submission
\usepackage{neurips_2026}
% the "default" option is equal to the "main" option, which is used for the Main Track with double-blind reviewing.
% 1. "main" option is used for the Main Track
%  \usepackage[main]{neurips_2026}
% 2. "position" option is used for the Position Paper Track
%  \usepackage[position]{neurips_2026}
% 3. "eandd" option is used for the Evaluations & Datasets Track
 % \usepackage[eandd]{neurips_2026}
% 4. "creativeai" option is used for the Creative AI Track
%  \usepackage[creativeai]{neurips_2026}
% 5. "sglblindworkshop" option is used for the Workshop with single-blind reviewing
 % \usepackage[sglblindworkshop]{neurips_2026}
% 6. "dblblindworkshop" option is used for the Workshop with double-blind reviewing
%  \usepackage[dblblindworkshop]{neurips_2026}
% 7. "education" option is used for the Education Track
%  \usepackage[education]{neurips_2026}


% After being accepted, the authors should add "final" behind the track to compile a camera-ready version.
% 1. Main Track
 % \usepackage[main, final]{neurips_2026}
% 2. Position Paper Track
%  \usepackage[position, final]{neurips_2026}
% 3. Evaluations & Datasets Track
 % \usepackage[eandd, final]{neurips_2026}
% 4. Creative AI Track
%  \usepackage[creativeai, final]{neurips_2026}
% 5. Workshop with single-blind reviewing
%  \usepackage[sglblindworkshop, final]{neurips_2026}
% 6. Workshop with double-blind reviewing
%  \usepackage[dblblindworkshop, final]{neurips_2026}
% Note. For the workshop paper template, both \title{} and \workshoptitle{} are required, with the former indicating the paper title shown in the title and the latter indicating the workshop title displayed in the footnote.
% For workshops (5., 6.), the authors should add the name of the workshop, "\workshoptitle" command is used to set the workshop title.
% \workshoptitle{WORKSHOP TITLE}
% 7. Education Track
% \usepackage[education, final]{neurips_2026}

% "preprint" option is used for arXiv or other preprint submissions
 % \usepackage[preprint]{neurips_2026}

% to avoid loading the natbib package, add option nonatbib:
%    \usepackage[nonatbib]{neurips_2026}

\usepackage[utf8]{inputenc} % allow utf-8 input
\usepackage[T1]{fontenc}    % use 8-bit T1 fonts
\usepackage{hyperref}       % hyperlinks
\usepackage{url}            % simple URL typesetting
\usepackage{booktabs}       % professional-quality tables
\usepackage{graphicx}       % figures
\usepackage{amsfonts}       % blackboard math symbols
\usepackage{nicefrac}       % compact symbols for 1/2, etc.
\usepackage{microtype}      % microtypography
\usepackage{xcolor}         % colors

% Note. For the workshop paper template, both \title{} and \workshoptitle{} are required, with the former indicating the paper title shown in the title and the latter indicating the workshop title displayed in the footnote. 
\title{A Shared Invertible Latent for Cross-Modal Medical Image Synthesis}


% The \author macro works with any number of authors. There are two commands
% used to separate the names and addresses of multiple authors: \And and \AND.
%
% Using \And between authors leaves it to LaTeX to determine where to break the
% lines. Using \AND forces a line break at that point. So, if LaTeX puts 3 of 4
% authors names on the first line, and the last on the second line, try using
% \AND instead of \And before the third author name.


\author{%
  David S.~Hippocampus\thanks{Use footnote for providing further information
    about author (webpage, alternative address)---\emph{not} for acknowledging
    funding agencies.} \\
  Department of Computer Science\\
  Cranberry-Lemon University\\
  Pittsburgh, PA 15213 \\
  \texttt{hippo@cs.cranberry-lemon.edu} \\
  % examples of more authors
  % \And
  % Coauthor \\
  % Affiliation \\
  % Address \\
  % \texttt{email} \\
  % \AND
  % Coauthor \\
  % Affiliation \\
  % Address \\
  % \texttt{email} \\
  % \And
  % Coauthor \\
  % Affiliation \\
  % Address \\
  % \texttt{email} \\
  % \And
  % Coauthor \\
  % Affiliation \\
  % Address \\
  % \texttt{email} \\
}


\begin{document}


\maketitle


\begin{abstract}
  Multi-modal medical imaging captures complementary views of the same
  pathology, motivating representations that are simultaneously
  \emph{modality-agnostic}, \emph{generative}, and \emph{directly usable}
  for downstream analysis. Replacing per-modality encoders and decoders with a
  single normalizing flow attains an exact shared latent but discards the
  convolutional inductive biases that medical images rely on, whereas training
  independent models per modality yields incompatible latent geometries. We
  present a hybrid that keeps a per-modality convolutional VAE and bridges the
  modality-specific latents through per-modality normalizing flows into one
  shared low-dimensional latent~$\mathbf{u}$, on which a Gaussian-mixture prior
  is placed. Because the flows are invertible, cross-modal translation is exact
  and routed entirely through~$\mathbf{u}$; because $\mathbf{u}$ is a native
  encoder output rather than a probed intermediate activation, it doubles as a
  representation for clustering and cross-modal alignment. On real ADNI
  T1$\leftrightarrow$FA (diffusion-derived fractional anisotropy), our model
  surpasses faithfully reproduced CFM, DDPM, DiT, and CycleGAN baselines on
  cross-modal SSIM and PSNR in both directions. Our central contribution is a
  controlled analysis of \emph{why} such a shared latent aligns: through a
  reconstruction-collapse ladder we show that cross-modal alignment is conferred
  neither by the mixture prior, the KL term, nor the flow likelihood, but by two
  necessary and jointly sufficient ingredients---a reconstruction-dominant
  objective and a low-dimensional bottleneck---isolating a failure mode common
  to multi-term latent-alignment models. We further find that the learned
  clusters are geometrically real and modality-consistent yet not aligned with
  diagnosis, a limit we trace to the evaluation cohort rather than the model,
  and we quantify the perception--distortion trade-off underlying the
  reconstruction metrics. Together these results give a compact, invertible
  recipe for cross-modal synthesis and a clear account of what does and does not
  drive latent alignment.
\end{abstract}


\section{Results}
\label{sec:results}

\begin{figure}[t]
  \centering
  \includegraphics[width=\linewidth]{figures/06_generation_process_all.png}
  \caption{Cross-modal generation process (T1$\to$FA), all methods; rows are each
  method's own trajectory, source to generated with GT at right. Our invertible
  shared-latent flow (top rows) transfers through a single $\mathbf{u}$.}
  \label{fig:genproc}
\end{figure}

\begin{table}[t]
  \centering
  \caption{Cross-modal synthesis (SSIM\,/\,PSNR).}
  \label{tab:gen}
  \small
  \begin{tabular}{lcccc}
    \toprule
    & \multicolumn{2}{c}{T1$\rightarrow$FA} & \multicolumn{2}{c}{FA$\rightarrow$T1} \\
    \cmidrule(lr){2-3}\cmidrule(lr){4-5}
    Method & SSIM & PSNR & SSIM & PSNR \\
    \midrule
    CycleGAN            & 0.680 & 17.18 & 0.743 & 17.94 \\
    DDPM                & 0.362 & 18.44 & 0.769 & 20.36 \\
    DiT                 & 0.665 & 14.32 & 0.757 & 18.28 \\
    MeanFlow            & 0.686 & 17.37 & 0.621 & 17.27 \\
    CFM                 & 0.732 & 18.46 & 0.802 & 19.99 \\
    \textbf{Ours}       & \textbf{0.764} & \textbf{19.95} & \textbf{0.859} & \textbf{22.05} \\
    \bottomrule
  \end{tabular}
\end{table}

\begin{table}[t]
  \centering
  \caption{Clustering ($K{=}4$). Sil$\uparrow$, DB$\downarrow$.}
  \label{tab:clust}
  \small
  \begin{tabular}{lccccc}
    \toprule
    Method & ACC & NMI & subj-CCA & Sil & DB \\
    \midrule
    Majority baseline & 0.395 & --    & --    & --    & --    \\
    naive VAE         & 0.442 & 0.112 & 0.047 & --    & --    \\
    CycleGAN          & 0.395 & 0.105 & --    & --    & --    \\
    DDPM              & 0.395 & 0.088 & --    & --    & --    \\
    MeanFlow          & 0.372 & 0.098 & --    & --    & --    \\
    DiT               & 0.442 & 0.091 & --    & --    & --    \\
    CFM               & 0.419 & 0.098 & --    & --    & --    \\
    Ours (default)    & 0.442 & 0.099 & 0.023 & 0.130 & 2.207 \\
    Ours (recon-opt)  & 0.349 & 0.059 & \textbf{0.628} & \textbf{0.181} & \textbf{1.760} \\
    \bottomrule
  \end{tabular}
\end{table}


\section{Ablations}
\label{sec:ablation}

\begin{table}[t]
  \centering
  \caption{Unified ablation. All variants on shared axes: bottleneck dim,
  reconstruction regime (mean vs.\ sum-dominant), self-recon SSIM (a collapse
  indicator: $0.68$=mean-brain), cross-modal SSIM, alignment CCA (chance $0.25$),
  and GMM cluster quality (silhouette Sil, Davies--Bouldin DB). All recon numbers
  are SSIM (PSNR in Table~\ref{tab:gen}). Alignment appears \emph{only} with
  sum-recon \emph{and} low dim; removing cross-recon breaks it; GMM strength moves
  Sil/DB, not synthesis.}
  \label{tab:abl}
  \small
  \begin{tabular}{llccccccc}
    \toprule
    & & & \multicolumn{3}{c}{SSIM} & align & \multicolumn{2}{c}{cluster} \\
    \cmidrule(lr){4-6}\cmidrule(lr){8-9}
    Variant & dim & recon & self & T1$\to$FA & FA$\to$T1 & CCA & Sil$\uparrow$ & DB$\downarrow$ \\
    \midrule
    \multicolumn{9}{l}{\emph{Collapse ladder (add one ingredient; balanced weights)}}\\
    bare-align                 & 100352 & mean & 0.68 & --    & --    & --   & --   & -- \\
    $+$cross/flow/GMM/KL       & $\le$100352 & mean & 0.68 & -- & -- & 0.00 & -- & -- \\
    flow-NLL $\times10/\times1000$ & 16 & mean & 0.68 & -- & -- & 0.00 & -- & -- \\
    A1/B1 (recon-dominant, hi-dim) & 100352 & \textbf{sum} & 0.89 & -- & -- & 0.00 & -- & -- \\
    \midrule
    \multicolumn{9}{l}{\emph{Loss deletion (sum-recon, dim 8)}}\\
    full                       & 8 & sum & 0.89 & 0.768 & 0.862 & 0.938 & -- & -- \\
    $-$GMM-NLL                 & 8 & sum & 0.89 & 0.769 & 0.862 & 0.952 & -- & -- \\
    $-$flow-NLL                & 8 & sum & 0.89 & 0.769 & 0.862 & 0.948 & -- & -- \\
    $-$KL                      & 8 & sum & 0.90 & 0.768 & 0.863 & 0.948 & -- & -- \\
    $-$paired-align            & 8 & sum & 0.89 & 0.767 & 0.863 & 0.951 & -- & -- \\
    $\mathbf{-}$\textbf{cross-recon} & 8 & sum & 0.90 & \textbf{0.724} & \textbf{0.797} & \textbf{0.802} & -- & -- \\
    \midrule
    \multicolumn{9}{l}{\emph{GMM flow-base, 3-stage (sum-recon; Sil(dx)$<0$ for all)}}\\
    v2 (weak GMM)              & 8  & sum & 0.89 & \textbf{0.770} & 0.863 & 0.94 & 0.314 & 1.211 \\
    v4 (strong S1-2, weak S3)  & 8  & sum & 0.89 & 0.768 & 0.863 & 0.94 & 0.305 & 1.137 \\
    v5 (16d, from scratch)     & 16 & sum & \textbf{0.90} & 0.765 & \textbf{0.864} & 0.94 & 0.264 & 1.371 \\
    v1 (strong GMM)            & 8  & sum & 0.88 & 0.763 & 0.854 & 0.93 & \textbf{0.585} & \textbf{0.579} \\
    \bottomrule
  \end{tabular}
\end{table}





\begin{figure}[t]
  \centering
  \includegraphics[width=\linewidth]{figures/16_3stage_u3d_gaussianity.png}
  \caption{Shared latent $\mathbf{u}$ (our run): 3-D layout coloured by GMM
  cluster (left) and per-cluster Gaussianity vs.\ $\chi(8)$ (right). Clusters are
  real but sub-Gaussian.}
  \label{fig:u3d}
\end{figure}

\begin{figure}[t]
  \centering
  \includegraphics[width=\linewidth]{figures/18_3stage_generation_process.png}
  \caption{Cross-modal generation emerges block-by-block through flow$^{-1}$
  (our run, both directions), staying on-manifold and progressive.}
  \label{fig:genblocks}
\end{figure}

\begin{figure}[t]
  \centering
  \includegraphics[width=\linewidth]{figures/17_3stage_recon.png}
  \caption{Bidirectional reconstruction (our run, 3 test slices, shared
  colourbar). Self-recon is bottleneck-limited; cross-modal detail is bounded by
  what the source modality carries.}
  \label{fig:recon}
\end{figure}

\begin{figure}[t]
  \centering
  \includegraphics[width=\linewidth]{figures/19-v4_3stage_flow_step_by_step.png}
  \caption{Full flow trace by disease case (run v4). Left: T1 morphs as
  $\mathrm{flow}_A$ pushes $z\!\to\!\mathbf{u}$; right: FA emerges as
  $\mathrm{flow}_B^{-1}$ pulls $\mathbf{u}\!\to\!z_B$.}
  \label{fig:flowtrace}
\end{figure}

% \appendix

% \section{Technical appendices and supplementary material}
% Technical appendices with additional results, figures, graphs, and proofs may be submitted with the paper submission before the full submission deadline (see above). You can upload a ZIP file for videos or code, but do not upload a separate PDF file for the appendix. There is no page limit for the technical appendices. 

% Note: Think of the appendix as ``optional reading'' for reviewers. The paper must be able to stand alone without the appendix; for example, adding critical experiments that support the main claims to an appendix is inappropriate. 

% %%%%%%%%%%%%%%%%%%%%%%%%%%%%%%%%%%%%%%%%%%%%%%%%%%%%%%%%%%%%

% \newpage
% \input{checklist.tex}


\end{document}